\section{Background}
\label{sec:background}

\subsection{Network and Submarine-Cable Infrastructure}
\label{sec:background_cross_layer}

Users reach services through domain names, CDN resources, replicas, and anycast
targets. Traceroute exposes the intervening IP hops and AS transitions. For
geographically separated regions, these logical paths depend on submarine
cables that join coastal landing stations to terrestrial backhaul and visible
ASes. Multiple cables can connect the same coastal regions, and one cable can
serve several landing locations.

The network-to-physical relationship is many-to-many. Different routers or AS
transitions may project onto one landing-region corridor, while a recurring
network transition may support several physical alternatives. We measure the
two distributions over the same observations: network diversity describes
observed logical transitions, and physical diversity describes their supported
corridor distribution.

Landing regions provide a stable geographic unit between exact landing
stations and country-wide cable counts. Exact stations preserve facility
detail but split nearby coastal access points that serve the same direction.
Country-level cable inventories combine distinct coastlines and international
directions. A bounded landing region retains coastal structure, and a corridor
identifies the paired entry-exit direction connecting two such regions.
Parallel cable systems can share one corridor while remaining distinguishable
in the underlying candidate records.

Countries provide a shared geographic and external-connectivity context.
Island, coastal, and landlocked geography shapes access to submarine
infrastructure; landing facilities, backhaul, gateways, exchanges, and operator
interconnection shape the paths available to users. Country grouping retains
variation among access networks and upstream providers within that context.

\subsection{Existing Cross-Layer Measurements}
\label{sec:background_existing_work}

Event-oriented studies relate new submarine connectivity to routing changes.
Bischof et al.\ characterized Cuba's connectivity before and after ALBA-1 and
later formulated the task of connecting Internet observations to the worldwide
cable mesh~\cite{bischof2015cuba,bischof2018untangling}. Fanou et al.\ measured
routing changes following new cable deployments~\cite{fanou2020unintended}.

User-oriented work measures submarine involvement in service paths. Liu et
al.\ combine resource discovery, RIPE Atlas traceroutes, geolocation, and
propagation feasibility to quantify how users reach Web resources through
submarine paths~\cite{liu2020outofsight}. Routing-oriented work connects cable
infrastructure to AS-level structure; Ye et al.\ infer AS relationships with
international cable and landing-station structures and analyze their
interaction with routing~\cite{ye2023impact}.

Exposure and concentration describe complementary aspects of service paths.
Exposure counts how frequently a service path enters a feasible submarine
candidate space. Concentration examines the structure of that space after
entry. A country may have high exposure and broad corridor support, or low
exposure with the exposed observations concentrated on a few directions.

Cross-layer cartography supplies geographic infrastructure representations.
Internet Atlas and InterTubes connect long-haul fiber facilities to observed
paths~\cite{durairajan2013atlas,durairajan2015intertubes}; iGDB integrates
facilities, fiber, and network entities~\cite{anderson2022igdb}; measurements of
intercontinental links identify layer-3 links and gateway routers
~\cite{carisimo2023hopaway}. Nautilus associates IP links with candidate cables
using traceroute, geolocation, latency, ownership, and cable metadata
~\cite{ramanathan2023nautilus}. Calypso combines cable layouts, network
relationships, latency, and country-level paths~\cite{wang2025calypso}, while
Xaminer applies cross-layer maps to configurable infrastructure analyses
~\cite{ramanathan2024xaminer}.

Service deployment determines which paths enter these mappings. Anycast
routing, replica placement, and interconnection select service instances and
their network paths~\cite{cicalese2015anycast,devries2017anycast}. We retain
this service-specific path view when projecting physical candidates and
comparing network and corridor distributions over identical atomic segments.

The paired construction aligns the statistical unit across layers. A network
transition and its corridor support originate from the same atomic segment,
so changes in concentration reflect the projection of a fixed observation
population. Country-service aggregation preserves deployment differences
among equivalent services and geographic differences among their users.

\subsection{Country-Service View of Cross-Layer Diversity}
\label{sec:background_motivation}

Country and service jointly define each observation population. Geographic and
interconnection environments shape source-side options; anycast, replica
selection, and routing choose among them. DNS Roots illustrate this
interaction. Functionally equivalent services operate independent anycast
deployments, so a country may reach different Roots through distinct
instances, AS transitions, and corridors. Distributed applications similarly
localize access in some countries and direct others toward remote
infrastructure.

Multi-target topology measurements provide a broader reference population.
Their dynamic destinations sample many network directions during the same
time window. Each service measurement instead retains the destinations
selected for that service. Comparing these families characterizes the
observed footprints of their target-selection and routing processes.

Our measurement separates three properties. \emph{Submarine exposure} is the
fraction of valid traceroutes containing an inter-region feasible corridor.
\emph{Corridor concentration} describes how candidate-bearing segment mass is
distributed across landing-region corridors. \emph{Cross-layer contraction}
compares that physical distribution with the network transitions formed from
the same segments. The framework next derives these quantities; Section
~\ref{sec:datasets} introduces the aligned path and infrastructure datasets.
